\documentclass[11pt,a4paper]{article}

\usepackage[utf8]{inputenc}
\usepackage[T1]{fontenc}
\usepackage[french]{babel}
\usepackage{geometry}
\usepackage{booktabs}
\usepackage{array}
\usepackage{multirow}
\usepackage{xcolor}
\usepackage{colortbl}
\usepackage{amsmath}
\usepackage{amssymb}
\usepackage{siunitx}
\usepackage{graphicx}
\usepackage{fancyhdr}
\usepackage{titlesec}
\usepackage{tcolorbox}
\usepackage{enumitem}

\geometry{margin=2.5cm}
\pagestyle{fancy}
\fancyhf{}
\fancyhead[L]{Simulation Geant4}
\fancyhead[R]{Comparaison Am-241 / Eu-152}
\fancyfoot[C]{\thepage}

\definecolor{headerblue}{RGB}{41, 128, 185}
\definecolor{alertred}{RGB}{192, 57, 43}
\definecolor{successgreen}{RGB}{39, 174, 96}
\definecolor{warningorange}{RGB}{230, 126, 34}
\definecolor{lightgray}{RGB}{245, 245, 245}

\titleformat{\section}{\Large\bfseries\color{headerblue}}{\thesection}{1em}{}
\titleformat{\subsection}{\large\bfseries}{\thesubsection}{1em}{}

\tcbuselibrary{skins,breakable}

\newtcolorbox{keypoint}[1][]{
    colback=headerblue!10,
    colframe=headerblue,
    fonttitle=\bfseries,
    title=#1,
    breakable
}

\newtcolorbox{warning}[1][]{
    colback=warningorange!10,
    colframe=warningorange,
    fonttitle=\bfseries,
    title=#1,
    breakable
}

\newtcolorbox{conclusion}[1][]{
    colback=successgreen!10,
    colframe=successgreen,
    fonttitle=\bfseries,
    title=#1,
    breakable
}

\newcolumntype{C}[1]{>{\centering\arraybackslash}p{#1}}

\begin{document}

\begin{center}
    {\LARGE\bfseries Analyse comparative des dépôts de dose}\\[0.4cm]
    {\Large Sources Am-241 et Eu-152 dans l'eau}\\[0.5cm]
    {\large Simulation Monte Carlo Geant4}\\[0.3cm]
    {\large 19 janvier 2026}
\end{center}

\vspace{0.5cm}

\tableofcontents

\newpage

%%%%%%%%%%%%%%%%%%%%%%%%%%%%%%%%%%%%%%%%%%%%%%%%%%%%%%%%%%%%%%%%%%%%%%%%%%%%%%%
\section{Contexte et problématique}
%%%%%%%%%%%%%%%%%%%%%%%%%%%%%%%%%%%%%%%%%%%%%%%%%%%%%%%%%%%%%%%%%%%%%%%%%%%%%%%

\subsection{Configuration géométrique}

La simulation modélise le dépôt de dose dans un volume d'eau de 3~mm d'épaisseur, divisé en deux zones :

\begin{center}
\begin{tabular}{lccc}
\toprule
\rowcolor{headerblue!20}
\textbf{Volume} & \textbf{Position (z)} & \textbf{Épaisseur} & \textbf{Comptabilisation} \\
\midrule
Water1 & 100--102 mm & 2 mm & \textcolor{alertred}{\textbf{Non comptée}} \\
\rowcolor{lightgray}
WaterRings (anneaux) & 102--103 mm & 1 mm & \textcolor{successgreen}{\textbf{Comptée}} \\
\bottomrule
\end{tabular}
\end{center}

\begin{warning}[Limitation actuelle]
La simulation ne comptabilise l'énergie déposée que dans les anneaux (1~mm), ignorant les 2~mm de Water1 situés en amont. Cette limitation sous-estime significativement le dépôt de dose pour les sources de basse énergie.
\end{warning}

\subsection{Observation initiale}

Pour une même activité de 42~kBq, les débits de dose mesurés dans l'anneau central (1~mm) sont étonnamment proches :

\begin{center}
\begin{tabular}{lccl}
\toprule
\rowcolor{headerblue!20}
\textbf{Source} & \textbf{Spectre principal} & \textbf{Débit (anneau 0)} & \textbf{Observation} \\
\midrule
Am-241 & 13--60 keV & $\sim 7{,}0~\mu$Gy/h & \multirow{2}{*}{Quasi-identiques !} \\
\rowcolor{lightgray}
Eu-152 & 122--1408 keV & $6{,}35~\mu$Gy/h & \\
\bottomrule
\end{tabular}
\end{center}

\textbf{Question :} Pourquoi des sources aux spectres si différents donnent-elles des doses similaires ?

%%%%%%%%%%%%%%%%%%%%%%%%%%%%%%%%%%%%%%%%%%%%%%%%%%%%%%%%%%%%%%%%%%%%%%%%%%%%%%%
\section{Analyse physique de la compensation}
%%%%%%%%%%%%%%%%%%%%%%%%%%%%%%%%%%%%%%%%%%%%%%%%%%%%%%%%%%%%%%%%%%%%%%%%%%%%%%%

\subsection{Facteurs influençant le dépôt de dose}

Le débit de dose dépend de plusieurs facteurs qui se compensent partiellement :

\begin{equation}
\dot{D} \propto A \times \bar{n}_\gamma \times \bar{E}_\gamma \times P_{\text{int}} \times f_{\text{dep}}
\end{equation}

où :
\begin{itemize}[noitemsep]
    \item $A$ : activité de la source (Bq)
    \item $\bar{n}_\gamma$ : nombre moyen de photons par désintégration
    \item $\bar{E}_\gamma$ : énergie moyenne des photons (keV)
    \item $P_{\text{int}}$ : probabilité d'interaction dans le volume
    \item $f_{\text{dep}}$ : fraction d'énergie effectivement déposée par interaction
\end{itemize}

\subsection{Comparaison des facteurs}

\begin{table}[h]
\centering
\begin{tabular}{lccC{2cm}}
\toprule
\rowcolor{headerblue!20}
\textbf{Facteur} & \textbf{Am-241} & \textbf{Eu-152} & \textbf{Ratio Am/Eu} \\
\midrule
Photons/désintégration ($\bar{n}_\gamma$) & 0,76 & $\sim 2{,}2$ & 0,35 \\
\rowcolor{lightgray}
Énergie moyenne ($\bar{E}_\gamma$) & $\sim 40$ keV & $\sim 400$ keV & 0,10 \\
Taux d'interaction ($P_{\text{int}}$) & $\sim 15\%$ & $\sim 2{,}5\%$ & 6,0 \\
\rowcolor{lightgray}
Fraction déposée ($f_{\text{dep}}$) & 100\% & $\sim 40\%$ & 2,5 \\
\midrule
\textbf{Produit global} & \textbf{1,0} & \textbf{$\sim 0{,}9$} & $\sim 1{,}1$ \\
\bottomrule
\end{tabular}
\caption{Facteurs de compensation entre Am-241 et Eu-152}
\end{table}

\subsection{Mécanismes d'interaction}

\begin{keypoint}[Effet photoélectrique (Am-241)]
\begin{itemize}[noitemsep]
    \item Dominant pour $E_\gamma < 100$~keV dans l'eau
    \item Section efficace $\sigma_{\text{pe}} \propto Z^5 / E^{3,5}$
    \item \textbf{100\% de l'énergie} du photon est transférée aux électrons
    \item Forte atténuation : $\mu \approx 1{-}3$~cm$^{-1}$ pour 15--60 keV
\end{itemize}
\end{keypoint}

\begin{keypoint}[Diffusion Compton (Eu-152)]
\begin{itemize}[noitemsep]
    \item Dominant pour $100~\text{keV} < E_\gamma < 5$~MeV
    \item Section efficace quasi-indépendante de $Z$
    \item \textbf{30--70\% de l'énergie} transférée (selon angle de diffusion)
    \item Faible atténuation : $\mu \approx 0{,}1{-}0{,}2$~cm$^{-1}$ pour 300--1000 keV
\end{itemize}
\end{keypoint}

%%%%%%%%%%%%%%%%%%%%%%%%%%%%%%%%%%%%%%%%%%%%%%%%%%%%%%%%%%%%%%%%%%%%%%%%%%%%%%%
\section{Impact de l'épaisseur considérée}
%%%%%%%%%%%%%%%%%%%%%%%%%%%%%%%%%%%%%%%%%%%%%%%%%%%%%%%%%%%%%%%%%%%%%%%%%%%%%%%

\subsection{Loi d'atténuation}

L'intensité transmise suit la loi de Beer-Lambert :
\begin{equation}
I(x) = I_0 \, e^{-\mu x}
\end{equation}

La fraction d'énergie absorbée entre 0 et $x$ est :
\begin{equation}
F(x) = 1 - e^{-\mu x}
\end{equation}

\subsection{Répartition de l'absorption selon l'énergie}

\begin{table}[h]
\centering
\begin{tabular}{ccccc}
\toprule
\rowcolor{headerblue!20}
\textbf{Énergie (keV)} & \textbf{$\mu$ (cm$^{-1}$)} & \textbf{Abs. 2 mm (\%)} & \textbf{Abs. 3 mm (\%)} & \textbf{Ratio 2mm/3mm} \\
\midrule
13,9 (X$_\text{L}\alpha$) & 2,5 & 39,3 & 52,8 & 0,74 \\
\rowcolor{lightgray}
17,0 (X$_\text{L}\beta$) & 1,5 & 25,9 & 36,2 & 0,72 \\
20,8 (X$_\text{L}\gamma$) & 0,9 & 16,5 & 23,7 & 0,70 \\
\rowcolor{lightgray}
59,5 ($\gamma$ principal) & 0,2 & 3,9 & 5,8 & 0,67 \\
\midrule
344 (Eu-152) & 0,11 & 2,2 & 3,2 & 0,67 \\
\bottomrule
\end{tabular}
\caption{Répartition de l'absorption dans Water1 (2 mm) et total (3 mm)}
\end{table}

\begin{warning}[Conséquence pour Am-241]
Pour les raies X de l'Am-241 (13--20 keV), \textbf{65 à 75\% de l'absorption totale} se produit dans les 2 premiers millimètres (Water1), qui ne sont pas comptabilisés dans la simulation actuelle !
\end{warning}

\subsection{Estimation du dépôt réel sur 3 mm}

En considérant que les anneaux (1 mm) ne représentent qu'environ 33\% de l'absorption totale pour l'Am-241 :

\begin{equation}
E_{\text{total}}^{\text{(3 mm)}} \approx \frac{E_{\text{anneaux}}^{\text{(1 mm)}}}{0{,}33} \approx 3 \times E_{\text{anneaux}}
\end{equation}

\begin{table}[h]
\centering
\begin{tabular}{lccc}
\toprule
\rowcolor{headerblue!20}
\textbf{Source} & \textbf{Débit 1 mm} & \textbf{Débit 3 mm (estimé)} & \textbf{Facteur} \\
\midrule
Am-241 & $7{,}0~\mu$Gy/h & $\mathbf{\sim 20~\mu}$\textbf{Gy/h} & $\times 2{,}9$ \\
\rowcolor{lightgray}
Eu-152 & $6{,}35~\mu$Gy/h & $\sim 7{-}8~\mu$Gy/h & $\times 1{,}1{-}1{,}2$ \\
\bottomrule
\end{tabular}
\caption{Comparaison des débits de dose selon l'épaisseur considérée}
\end{table}

%%%%%%%%%%%%%%%%%%%%%%%%%%%%%%%%%%%%%%%%%%%%%%%%%%%%%%%%%%%%%%%%%%%%%%%%%%%%%%%
\section{Synthèse comparative}
%%%%%%%%%%%%%%%%%%%%%%%%%%%%%%%%%%%%%%%%%%%%%%%%%%%%%%%%%%%%%%%%%%%%%%%%%%%%%%%

\subsection{Tableau récapitulatif}

\begin{table}[h]
\centering
\renewcommand{\arraystretch}{1.3}
\begin{tabular}{lcc}
\toprule
\rowcolor{headerblue!20}
\textbf{Paramètre} & \textbf{Am-241} & \textbf{Eu-152} \\
\midrule
Activité & \multicolumn{2}{c}{42 kBq} \\
\rowcolor{lightgray}
Spectre principal & 13--60 keV & 122--1408 keV \\
Interaction dominante & Photoélectrique & Compton \\
\rowcolor{lightgray}
$\mu$ moyen (cm$^{-1}$) & $\sim 1{-}2$ & $\sim 0{,}1$ \\
\midrule
\multicolumn{3}{c}{\cellcolor{headerblue!10}\textbf{Débits de dose}} \\
\midrule
Anneau central (1 mm) & $7{,}0~\mu$Gy/h & $6{,}35~\mu$Gy/h \\
\rowcolor{lightgray}
Région centrale (3 mm) & $\mathbf{\sim 20~\mu}$\textbf{Gy/h} & $\sim 7{-}8~\mu$Gy/h \\
\textbf{Ratio Am-241 / Eu-152} & \multicolumn{2}{c}{\textbf{$\sim 2{,}7$}} \\
\midrule
\multicolumn{3}{c}{\cellcolor{headerblue!10}\textbf{Temps pour 20 cGy (région centrale)}} \\
\midrule
Calcul & $\dfrac{200\,000~\mu\text{Gy}}{20~\mu\text{Gy/h}}$ & $\dfrac{200\,000~\mu\text{Gy}}{7{,}5~\mu\text{Gy/h}}$ \\[0.3cm]
\rowcolor{lightgray}
Résultat & $\mathbf{\sim 417}$ \textbf{jours} & $\sim 1\,111$ jours \\
& $\approx$ \textbf{1,1 an} & $\approx$ 3 ans \\
\bottomrule
\end{tabular}
\caption{Synthèse comparative Am-241 vs Eu-152}
\end{table}

\subsection{Explication de la compensation apparente (1 mm)}

Sur le dernier millimètre uniquement, les débits sont proches car :

\begin{center}
\renewcommand{\arraystretch}{1.4}
\begin{tabular}{rcl}
\textbf{Am-241} & : & $\underbrace{\text{Moins de } \gamma}_{\times 0{,}35} \times \underbrace{\text{Forte absorption}}_{\times 6} \times \underbrace{100\%~E~\text{déposée}}_{\times 2{,}5} \times \underbrace{E~\text{faible}}_{\times 0{,}1}$ \\[0.2cm]
& & Produit $\approx 0{,}35 \times 6 \times 2{,}5 \times 0{,}1 \approx 0{,}5$ \\[0.3cm]
\textbf{Eu-152} & : & $\underbrace{\text{Plus de } \gamma}_{\times 1} \times \underbrace{\text{Faible interact.}}_{\times 1} \times \underbrace{40\%~E~\text{déposée}}_{\times 1} \times \underbrace{E~\text{élevée}}_{\times 1}$ \\[0.2cm]
& & Produit $\approx 1 \times 1 \times 1 \times 1 = 1$ \\
\end{tabular}
\end{center}

Le ratio $\sim 0{,}5$ pour Am-241 est compensé par le fait que la majorité de l'absorption s'est déjà produite en amont (Water1), ne laissant qu'une fraction résiduelle dans les anneaux.

%%%%%%%%%%%%%%%%%%%%%%%%%%%%%%%%%%%%%%%%%%%%%%%%%%%%%%%%%%%%%%%%%%%%%%%%%%%%%%%
\section{Conclusions}
%%%%%%%%%%%%%%%%%%%%%%%%%%%%%%%%%%%%%%%%%%%%%%%%%%%%%%%%%%%%%%%%%%%%%%%%%%%%%%%

\begin{conclusion}[Résultats principaux]
\begin{enumerate}
    \item \textbf{La proximité des débits sur 1 mm est un artefact} de la géométrie de mesure, non représentatif du dépôt total.
    
    \item \textbf{Sur l'épaisseur complète de 3 mm}, l'Am-241 dépose \textbf{2,5 à 3 fois plus de dose} que l'Eu-152, conformément à l'intuition physique.
    
    \item \textbf{Le temps pour atteindre 20 cGy} dans la région centrale est :
    \begin{itemize}
        \item Am-241 : $\sim 1{,}1$ an (contre 3,3 ans si on ne compte que 1 mm)
        \item Eu-152 : $\sim 3$ ans
    \end{itemize}
    
    \item \textbf{Recommandation} : Modifier la simulation pour comptabiliser l'énergie déposée dans Water1, ou fusionner Water1 et les anneaux en un volume unique de 3 mm.
\end{enumerate}
\end{conclusion}

\vspace{0.5cm}

\begin{keypoint}[Cohérence physique validée]
Les résultats sont cohérents avec la physique des interactions photon-matière :
\begin{itemize}
    \item Fort pouvoir d'arrêt des photons de basse énergie (photoélectrique)
    \item Absorption concentrée dans les premiers millimètres pour l'Am-241
    \item Absorption quasi-uniforme pour l'Eu-152 (Compton)
\end{itemize}
\end{keypoint}

\vspace{1cm}

\begin{center}
\begin{tabular}{|c|c|c|}
\hline
\rowcolor{headerblue!30}
& \textbf{Am-241} & \textbf{Eu-152} \\
\hline
\textbf{Débit dose (3 mm)} & $\sim 20~\mu$Gy/h & $\sim 7{,}5~\mu$Gy/h \\
\hline
\textbf{Temps pour 20 cGy} & \cellcolor{successgreen!30}\textbf{1,1 an} & 3 ans \\
\hline
\end{tabular}
\end{center}

\vspace{0.5cm}

\begin{flushright}
\textit{Document généré le 19 janvier 2026}
\end{flushright}

\end{document}
